\section{实验评估}
\label{sec:experiments}
我们回答三个研究问题。RQ 1 考察 \sys{} 是否<……>(贡献 1)。RQ 2 考察<……>。RQ 3 考察<……>。实验覆盖 \num{families} 个方法族、\num{baselines} 个基线和 \num{datasets} 个数据集。

\subsection{实验设置}
\label{sec:setup}
(1)\textbf{数据集。}表~\ref{tab:datasets}列出了各数据集的任务、错误类型与规模。<数据从哪来、错误从哪来>。
(2)\textbf{基线。}我们与表~\ref{tab:related}中的全部方法比较。<所有方法从什么数据学习>。<\sys{} 有何不同>。
(3)\textbf{指标。}我们报告\textbf{召回率}(\textbf{R})，$R = \tfrac{|\text{找到的重复}|}{|\text{全部重复}|}$，以及<……>。

\begin{table}[t]
\centering
\caption{各数据集的任务、错误类型与规模。}
\label{tab:datasets}
\footnotesize
\begin{tabular}{@{}r|c|c|c@{}}
\toprule
Dataset & Task & Error types & Records \\
\midrule
\textbf{<D1>}~\cite{placeholder2024} & DD & T/M & <n> \\
\textbf{<D2>}~\cite{placeholder2024} & DD & T/F & <n> \\
\bottomrule
\end{tabular}
\end{table}

\subsection{实验结果与发现}
\label{sec:findings}

\begin{table}[t]
\centering
\caption{相同标注预算下各方法的召回率。}
\label{tab:main}
\footnotesize
\setlength{\tabcolsep}{3.6pt}
\renewcommand{\arraystretch}{0.9}
\begin{tabular}{@{}r@{\hspace{4pt}}r|ccc@{}}
\toprule
Rank & Method & D1$\uparrow$ & D2$\uparrow$ & D3$\uparrow$ \\
\midrule
%%BEGIN-MAINROWS
\grp{5}{Rule-based} \\
4 & R-Match & 0.612 & 0.574 & 0.655 \\
5 & R-Block & 0.598 & 0.603 & 0.621 \\
\grp{5}{Learning-based} \\
3 & L-Forest & 0.844 & 0.792 & \secondbest{0.861} \\
2 & L-Boost & \secondbest{0.851} & \secondbest{0.805} & 0.833 \\
\midrule
\oursrow
1 & \textbf{\sys} & \best{0.912} & \best{0.866} & \best{0.890} \\
%%END-MAINROWS
\bottomrule
\end{tabular}
\end{table}

\begin{finding}
\label{find:main}
\sys{} <带有它所胜出的比较的强论断>。
\end{finding}
\sys{} 比最强的基线高 \num{gain over best pct}\%(表~\ref{tab:main})。<机理，对应引言的观察 (1)>。<边界情形及其实测原因>。

\begin{finding}
\label{find:cost}
<关于代价或规模的论断>。
\end{finding}
\sys{} 在最大的表上用时 \num{runtime minutes} 分钟。<哪一步占主要开销，随规模如何增长>。
